\section{The Privacy Externality Game}
\label{sec:game}

We formalise the strategic situation faced by data contributors in
sampling-based iDP as a finite normal-form game.

\subsection{Players, actions and mechanism}

\begin{definition}[Privacy Externality Game]
\label{def:peg}
Fix a common $\delta>0$, an expected batch size $E_b$, an iteration count $I$,
and a set of $n$ \emph{groups} with sizes $|G_1|,\dots,|G_n|$. The Privacy
Externality Game (PEG) is the tuple $\langle N, \Sigma, U\rangle$ where
\begin{itemize}
\item $N = \{1,\dots,n\}$ is the set of players. Each player $i$ controls
group $i$.
\item Each player's action set is a finite grid
$\Sigma_i = \mathcal{A} \subset \R_{>0}$ of admissible
$\eps$-values, $\eps_i^{\min} \leq \eps_i \leq \eps_i^{\max}$.
\item Given the action profile $\boldsymbol{\eps} \in \mathcal{A}^n$, the
mechanism solution $(\sigma(\boldsymbol{\eps}), p_1(\boldsymbol{\eps}),\dots,
p_n(\boldsymbol{\eps}))$ is determined by Definition~\ref{def:sampling-idp},
provided the solution exists; if it does not, we say $\boldsymbol{\eps}$ is
\emph{infeasible}.
\item Each player's utility $u_i(\boldsymbol{\eps})$ depends on
$\sigma(\boldsymbol{\eps})$, on the realised MIA advantage
$A_i(\boldsymbol{\eps}) := \Adv\bigl(\mathcal{M}_{p_i,\sigma,I}\bigr)$, and on
the player's own action $\eps_i$ (specified per regime below).
\end{itemize}
\end{definition}

\begin{assumption}[Feasibility]
\label{ass:feasibility}
The action grid $\mathcal{A}^n$ is restricted to its feasible subset
$\mathcal{F} \subseteq \mathcal{A}^n$ where the mechanism solution exists.
Infeasible profiles are removed from a player's action set after the others'
moves are observed.
\end{assumption}

\subsection{Vulnerability function and externality}

\begin{definition}[Vulnerability function]
\label{def:vulnerability}
The \emph{vulnerability function} of player $i$ is
\[
A_i\colon \mathcal{F} \to [0, 1],\quad \boldsymbol{\eps} \mapsto
\Adv\bigl(\mathcal{M}_{p_i(\boldsymbol{\eps}),\,\sigma(\boldsymbol{\eps}),\,I}\bigr).
\]
\end{definition}

In Pej\'o et al.'s \emph{Together or Alone}~\cite{pejo2019together}, the
analogous function $\Phi_n(p_1,p_2)$ measured \emph{accuracy loss} arising
from a player's own privacy choice. In our setting, $A_i(\boldsymbol{\eps})$
measures \emph{realised privacy risk} arising from \emph{every player's} choice.
This is the privacy externality.

\subsection{Awareness regimes}

We separate two contracting situations. Both occur in practice and the
distinction is critical to interpreting the equilibrium.

\paragraph{Naive regime.}
The player perceives the cost of disclosure as her \emph{nominal} budget
$\eps_i$:
\begin{equation}
u_i^{\textrm{naive}}(\boldsymbol{\eps}) \;=\; B_i\, V(\sigma(\boldsymbol{\eps}))
- C_i\, \frac{\eps_i}{\eps_i^{\max}}.
\label{eq:utility-naive}
\end{equation}

\paragraph{Aware regime.}
The player perceives the cost of disclosure as the \emph{realised} MIA
advantage:
\begin{equation}
u_i^{\textrm{aware}}(\boldsymbol{\eps}) \;=\; B_i\, V(\sigma(\boldsymbol{\eps}))
- C_i\, A_i(\boldsymbol{\eps}).
\label{eq:utility-aware}
\end{equation}

Here $V$ is a model-utility proxy that is monotone-decreasing in $\sigma$ (we
use $V(\sigma) = 1/(1+\sigma)$ throughout) and $B_i, C_i > 0$ are player-
specific weights mirroring the $B$ and $C$ in
Pej\'o et al.~\cite{pejo2019together}. The naive regime models a player who
has been given the formal $(\eps_i,\delta)$ guarantee and prices it
verbatim; the aware regime models a player who knows the realised vulnerability
of her data after the others have played.

\begin{remark}
A naive player who knows only her contracted $(\eps_i,\delta)$ cannot detect
that her actual risk has increased through her peers' choices. The realised
MIA advantage $A_i(\boldsymbol{\eps})$ may exceed her contractually expected
vulnerability without any formal guarantee being broken. We make this point
formal in Lemma~\ref{lem:externality}.
\end{remark}

\subsection{Coalition variant}

\begin{definition}[$k$-coalition Stackelberg PEG]
\label{def:coalition}
Fix a target index $t \in N$ and a colluding subset $C \subset N \setminus
\{t\}$ with $|C| = k$. The $k$-coalition Stackelberg PEG is the two-stage game
in which:
\begin{enumerate}
\item Stage~1 (leader): the coalition jointly commits to an action vector
$\boldsymbol{\eps}_C \in \mathcal{A}^k$ to maximise
\(
\lambda \cdot A_t(\boldsymbol{\eps}) + \sum_{i \in C} u_i^{\textrm{aware}}(\boldsymbol{\eps}),
\)
where $\lambda \geq 0$ controls the coalition's appetite for harming~$t$ vs.\
maintaining their own utility.
\item Stage~2 (followers): the non-coalition players $\{t\} \cup (N \setminus
(C \cup \{t\}))$ best-respond.
\end{enumerate}
\end{definition}

\subsection{Equilibrium concepts}

We will use two equilibrium concepts. The \emph{pure Nash equilibrium} of the
PEG is an $\boldsymbol{\eps}^* \in \mathcal{F}$ such that
$u_i(\eps_i^*, \boldsymbol{\eps}_{-i}^*) \geq u_i(\eps_i, \boldsymbol{\eps}_{-i}^*)$
for every player $i$ and every $\eps_i \in \mathcal{A}$ such that
$(\eps_i, \boldsymbol{\eps}_{-i}^*) \in \mathcal{F}$. Where no pure NE exists,
we characterise the empirical mixed equilibrium reached by fictitious
play~\cite{monderer1996potential}.

\begin{definition}[Price of Excess Vulnerability]
\label{def:poev}
Let $\boldsymbol{\eps}^{(0)}$ be a baseline profile (in our experiments,
$\boldsymbol{\eps}^{(0)} = (\bar\eps,\dots,\bar\eps)$ for the mean of the
admissible grid). For an equilibrium $\boldsymbol{\eps}^*$, the
\emph{Price of Excess Vulnerability} is
\begin{equation}
\PoEV(\boldsymbol{\eps}^*)\;=\;
\frac{\sum_{i\in N} A_i(\boldsymbol{\eps}^*) - \sum_{i\in N} A_i(\boldsymbol{\eps}^{(0)})}
{\sum_{i\in N} A_i(\boldsymbol{\eps}^{(0)})}.
\label{eq:poev}
\end{equation}
\end{definition}
$\PoEV \geq 0$ measures the population-level inflation of MIA advantage at the
equilibrium relative to the uniform baseline. It is the privacy analogue of
Pej\'o et al.'s Price of Privacy, expressed in vulnerability rather than
accuracy loss.
